%=========================================================
%               CLASSE DU DOCUMENT
%=========================================================
\documentclass[12pt,a4paper]{article}

%=========================================================
%         ENCODAGE ET POLICES DE CARACTÈRES
%=========================================================
\usepackage[utf8]{inputenc}    % Encodage UTF-8 (inutile avec XeLaTeX/LuaLaTeX)
\usepackage[T1]{fontenc}       % Meilleure gestion des caractères accentués
\usepackage[french]{babel} % <--- Ce package fournit \og et \fg{}
\usepackage{lmodern}           % Police Latin Modern
\usepackage{times}             % Police Times (à supprimer si tu préfères Latin Modern)

%=========================================================
%                MATHÉMATIQUES
%=========================================================
\usepackage{amsmath}           % Environnements mathématiques avancés
\usepackage{amssymb}           % Symboles mathématiques supplémentaires
\usepackage{amsfonts}          % Polices mathématiques
\usepackage{mathrsfs}          % Police calligraphique (\mathscr)
\usepackage{tkz-tab}           % Tableaux de signes et de variations
\everymath{\displaystyle}      % Affiche toutes les formules en style "display"

%=========================================================
%               MISE EN PAGE
%=========================================================
\usepackage[
    left=1cm,
    right=1cm,
    top=1cm,
    bottom=1.7cm
]{geometry}                    % Marges du document

\usepackage{multicol}          % Texte sur plusieurs colonnes
\usepackage{float}             % Placement précis des figures et tableaux

%=========================================================
%            TABLEAUX ET LISTES
%=========================================================
\usepackage{array}             % Colonnes de tableaux personnalisées
\usepackage{multirow}          % Fusion de lignes dans les tableaux
\usepackage{enumitem}          % Personnalisation des listes

\usepackage{tikz}
\usetikzlibrary{angles,quotes,arrows.meta}
%=========================================================
%            COULEURS ET BOÎTES
%=========================================================
\usepackage[most]{tcolorbox}   % Boîtes colorées très personnalisables
\tcbuselibrary{skins,hooks}    % Bibliothèques supplémentaires de tcolorbox

\usepackage{mdframed}          % Cadres autour de paragraphes
\usepackage{varwidth}          % Largeur automatique des boîtes

%=========================================================
%                DESSINS ET GRAPHIQUES
%=========================================================
\usepackage{tikz}              % Dessins vectoriels
\usetikzlibrary{calc}
\usetikzlibrary{
    arrows,
    shapes.geometric,
    fit,
    patterns
}

\usepackage{pgfplots}          % Tracés de fonctions et graphiques
\pgfplotsset{compat=1.15}

%=========================================================
%              LIENS HYPERTEXTES
%=========================================================
\usepackage[
    colorlinks=true,
    linkcolor=blue,
    urlcolor=blue,
    citecolor=blue
]{hyperref}

%=========================================================
%          ÉLÉMENTS EN ARRIÈRE-PLAN
%=========================================================
\usepackage{eso-pic}           % Ajout de filigranes ou d'éléments sur chaque page

%=========================================================
%              SYMBOLES DIVERS
%=========================================================
\usepackage{pifont}            % Symboles (✓, ✗, etc.)

\DeclareUnicodeCharacter{2460}{\text{\ding{172}}} % Pour ①
\DeclareUnicodeCharacter{2461}{\text{\ding{173}}} % Pour ②
%=========================================================
%      PERSONNALISATION DES LISTES NUMÉROTÉES
%=========================================================

\newcommand\circitem[1]{%
\tikz[baseline=(char.base)]{
\node[circle,draw=gray,fill=red!55,
minimum size=1.2em,inner sep=0] (char) {#1};}}

\newcommand\boxitem[1]{%
\tikz[baseline=(char.base)]{
\node[fill=cyan,
minimum size=1.2em,inner sep=0] (char) {#1};}}

\setlist[enumerate,1]{label=\protect\circitem{\arabic*}}
\setlist[enumerate,2]{label=\protect\boxitem{\alph*}}

%=========================================================
%            COMMANDES PERSONNALISÉES
%=========================================================

% Soulignement coloré
\newcommand{\myul}[2][black]{%
\setulcolor{#1}\ul{#2}\setulcolor{black}}

% Tabulation horizontale
\newcommand\tab[1][1cm]{\hspace*{#1}}

%=========================================================
%             COMPTEURS PERSONNALISÉS
%=========================================================

%---------- Exemple ----------
\newcounter{exemple}

\newcommand{\exemple}{%
\refstepcounter{exemple}
\textbf{\textcolor{orange}{Exemple \theexemple : }}
\ignorespaces}

%---------- Solution ----------
\newcounter{solution}

\newcommand{\solution}{%
\refstepcounter{solution}
\textbf{\textcolor{orange}{Solution \thesolution : }}
\ignorespaces}

%---------- Exercices ----------
\definecolor{myorange}{rgb}{1.0,0.8,0.0}

\newcounter{exerciceapp}

\newcommand{\exerciceapp}{%
\refstepcounter{exerciceapp}
\textbf{\textcolor{myorange}
{Exercice d'application \theexerciceapp : }}
\ignorespaces}

%---------- Corrections ----------
\newcounter{correction}

\newcommand{\correction}{%
\refstepcounter{correction}
\textbf{\textcolor{myorange}
{Correction \thecorrection : }}
\ignorespaces}

%---------- Remarques ----------
\definecolor{myorange1}{rgb}{1.0,1.0,0.0}
\usepackage{enumitem}

\definecolor{myred}{RGB}{180, 40, 40}
\definecolor{myblue}{RGB}{20, 50, 120}
\newcounter{remarque}

\newcommand{\remarque}{%
\refstepcounter{remarque}
\textbf{\textcolor{myorange1}
{Remarque \theremarque : }}
\ignorespaces}

%=========================================================
%                 FILIGRANE
%=========================================================

\AddToShipoutPicture{
    \AtTextCenter{
        \makebox[0pt]{
            \rotatebox{45}{
                \textcolor[gray]{0.6}{
                    \fontsize{5cm}{5cm}\selectfont PGB
                }
            }
        }
    }
}

\begin{document}
% En-tête personnalisée
\begin{center}
    \Large\textbf{\underline{\textcolor{red}{Limite et Continuité}}}\\[-0.1cm]
    \normalsize\textbf{Prof : M. BA} \hfill \textbf{Classe : 1erS2}\\[-0.1cm]
    \textbf{Année scolaire : 2025 -- 2026}
\end{center}

\section*{Objectifs spécifiques}
\begin{itemize}[label=$\checkmark$]
    \item Calculer des limites
\end{itemize}

\section*{Prérequis}
\begin{itemize}[label=$\checkmark$]
    \item Limites (vues en 1ère)
\end{itemize}

\section*{Supports didactiques}
\begin{itemize}[label=$\checkmark$]
    \item Cours limite et continuité de Bousso TS2;
    \item Mon cours de limite et continuité de TS2;
    \item CIAM 1ère L;
    \item Dimathème Terminale A1 et A2.
\end{itemize}

\section*{Plan du chapitre}
\begin{enumerate}[label=\Roman*.]
    \item \textbf{Limites de fonctions usuelles}
    \begin{enumerate}
        \item Théorème 1
        \item Théorème 2
        \item Théorème 3
    \end{enumerate}
    \item \textbf{Opérations sur les limites}
    \begin{enumerate}
        \item Limite d'une somme
        \item Limite d'un produit
        \begin{enumerate}
            \item Limite à l'infini d'un polynôme
        \end{enumerate}
        \item Limite d'un quotient
        \begin{enumerate}
            \item Limite à l'infini d'une fraction rationnelle
            \item Exemples de calculs de limites à gauche et à droite en un réel
        \end{enumerate}
    \end{enumerate}
\end{enumerate}

\newpage
\section*{\color{myred} I. Limites de fonctions usuelles }
\subsection*{\color{myred} 1. Théorème 1 }

Soit \( a \) un réel constant ou \( a = \pm \infty \) et \( c \) est un réel constant. On a alors :
$
\lim_{x \to a} c = c
$

\paragraph{Exemples :}
$
\lim\limits_{x \to 2} 3 = 3 \quad \text{et} \quad \lim\limits_{x \to +\infty} -5 = -5
$

\subsection*{\color{myred} 2. Théorème 2 }
Soit \( n \) un entier naturel non nul, on a alors :

\begin{enumerate}[label=\roman*.]
    \item 
    $
    \lim\limits_{x \to +\infty} x^n = +\infty
    $
    \paragraph{Exemples :}
    $    \lim\limits_{x \to +\infty} x^2 = +\infty \quad \text{et} \quad \lim\limits_{x \to +\infty} x^5 = +\infty$

    \item 
     $\lim\limits_{x \to -\infty} x^n = 
    \begin{cases}
        +\infty & \text{si } n \text{ est pair}, \\
        -\infty & \text{si } n \text{ est impair}.
    \end{cases}$
    \paragraph{Exemples :}
$    \lim\limits_{x \to -\infty} x^2 = +\infty \quad \text{et} \quad \lim_{x \to -\infty} x^7 = -\infty$

    \item 
$    \lim\limits_{x \to +\infty} \frac{1}{x^n} = 0$
    \paragraph{Exemples :}
$    \lim\limits_{x \to +\infty} \frac{1}{x^2} = 0 \quad \text{et} \quad \lim\limits_{x \to +\infty} \frac{1}{x^4} = 0$
\end{enumerate}

\subsection*{\color{myred} 3. Théorème 3 }
Si \( P(x) \) est une fonction polynôme et \( a \) est un réel constant alors :
$\lim\limits_{x \to a} P(x) = P(a)$

\paragraph{Exemple :}
$\lim\limits_{x \to -1} (x^3 - 2x^2 + 3x + 7) = 1$

\section*{\color{myred}II. Opérations sur les limites}

Dans les tableaux suivants, \( f(x) \) et \( g(x) \) sont des fonctions, \( a \) est un réel constant ou bien \( a = \pm\infty \), \( L \) et \( L' \) sont des réels constants.

\subsection*{\color{myred}1. Limite d'une somme}

\[
\begin{array}{|c|c|c|c|c|c|c|}
\hline
\lim_{x \to a} f(x) & L & L & L & +\infty & -\infty & +\infty \\
\hline
\lim_{x \to a} g(x) & L' & +\infty & -\infty & +\infty & -\infty & -\infty \\
\hline
\lim_{x \to a} (f(x) + g(x)) & L + L' & +\infty & -\infty & +\infty & -\infty & \text{Forme indéterminée} \\
\hline
\end{array}
\]

\paragraph{Exemples :}
\begin{itemize}
    \item  $\lim_{x \to +\infty} \left(x^3 + \frac{1}{x^2}\right) = +\infty$

    \item $    \lim_{x \to -\infty} \left(x^5 + \frac{1}{x^4}\right) = -\infty$

\end{itemize}

\subsection*{\color{myred}2. Limite d'un produit}

\subsubsection*{\color{myred}a. Limites de \( \alpha \times f(x) \)}

\( \alpha \) est un réel constant, on a le tableau suivant :

\[
\begin{array}{|c|c|c|c|}
\hline
\lim_{x \to a} f(x) & L & +\infty & -\infty \\
\hline
\lim_{x \to a} (\alpha \times f(x)) & \alpha \times L & 
\begin{array}{c}
+\infty \text{ si } \alpha > 0 \\
-\infty \text{ si } \alpha < 0
\end{array} & 
\begin{array}{c}
-\infty \text{ si } \alpha > 0 \\
+\infty \text{ si } \alpha < 0
\end{array} \\
\hline
\end{array}
\]

\paragraph{Exemples :}
\begin{itemize}
    \item Calculons \( \lim\limits_{x \to +\infty} 2x^2\) : \\
    \(\lim\limits_{x \to +\infty} x^2 = +\infty\) et \(2 > 0\) donc 
$ \lim\limits_{x \to +\infty} 2x^2 = +\infty$

    \item Calculons \(\lim\limits_{x \to +\infty} -2x^2\) : \\
    \( \lim\limits_{x \to +\infty} x^2 = +\infty\) et \(-2 < 0\) donc 
$ \lim\limits_{x \to +\infty} -2x^2 = -\infty$
\end{itemize}

\subsubsection*{\color{myred}b. Théorème}

La limite à l'infini d'un polynôme est égale à la limite à l'infini de son monôme de plus haut degré.

\paragraph{Exemples :}
\[
\lim_{x \to +\infty} (-3x^4 - 2x + 7) = \lim_{x \to +\infty} (-3x^4) = -\infty
\]
et
\[
\lim_{x \to +\infty} (2x^3 - 5x + 1) = \lim_{x \to +\infty} (2x^3) = +\infty
\]

\subsubsection*{\color{myred}c. Limites de \( f(x) \times g(x) \)}

\[
\begin{array}{|c|c|c|c|c|c|c|c|c|c|}
\hline
\lim_{x \to a} f(x) & 0 & >0 & >0 & <0 & <0 & +\infty & +\infty & -\infty & 0 \\
\hline
\lim_{x \to a} g(x) & L' & +\infty & -\infty & +\infty & -\infty & +\infty & -\infty & -\infty & \infty \\
\hline
\lim_{x \to a} (f(x) \times g(x)) & 0 & +\infty & -\infty & -\infty & +\infty & +\infty & -\infty & +\infty & \text{F.I.} \\
\hline
\end{array}
\]

\subsection*{\color{myred}3. Limites d'un quotient}

\subsubsection*{\color{myred}a. Cas où la limite du dénominateur est non nulle}

\[
\begin{array}{|c|c|c|c|c|c|c|c|}
\hline
\lim_{x \to a} f(x) & L & L & +\infty & +\infty & -\infty & -\infty & \infty \\
\hline
\lim_{x \to a} g(x) & L' \neq 0 & \infty & L' > 0 & L' < 0 & L' > 0 & L' < 0 & \infty \\
\hline
\lim_{x \to a} \frac{f(x)}{g(x)} & \frac{L}{L'} & 0 & +\infty & -\infty & -\infty & +\infty & \text{F.I.} \\
\hline
\end{array}
\]

\subsubsection*{Limite à l'infini d'une fraction rationnelle}

La limite à l'infini d'une fraction rationnelle est égale à la limite à l'infini du quotient du monôme de plus haut degré de son numérateur par le monôme de plus haut degré de son dénominateur.


\paragraph{Exemple :}
\[
\lim_{x \to +\infty} \frac{x^3 + x + 6}{x^2 + x - 1} = \lim_{x \to +\infty} \frac{x^3}{x^2} = \lim_{x \to +\infty} x = +\infty
\]

\subsubsection*{\color{myred}b. Cas où la limite du dénominateur est nulle}

\[
\begin{array}{|c|c|c|c|}
\hline
\lim_{x \to a} f(x) & L \neq 0 & \infty & 0 \\
\hline
\lim_{x \to a} g(x) & 0 & 0 & 0 \\
\hline
\lim_{x \to a} \frac{f(x)}{g(x)} & \infty & \infty & \text{F.I.} \\
\hline
\end{array}
\]

Dans les deux premiers cas, pour savoir de quel infini il s'agit, on doit étudier le signe du dénominateur.

\paragraph{Exemples :}

\paragraph{Exemple 1}
Étudions 
\[
\lim_{x \to 2} \frac{3x - 2}{2x - 4}
\]

\paragraph{Exemple 2}
Étudions 
\[
\lim_{x \to 1} \frac{x^2 + x + 6}{2x^2 + x - 1}
\]
\end{document}